\mytasksection{Задача 3. Подготовить систему для внедрения в организации}

Подготовка \emph{Onboarding System} (системы адаптации новых сотрудников) включала создание контейнер-образов для frontend и backend, настройку синхронизации с кадровой системой RIMS и подготовку конфигураций для корпоративного внедрения.

\subsubsection*{Контейнеризация компонентов}

\textbf{Backend контейнер (ASP.NET Core 9):}

Dockerfile реализует двухэтапную сборку: на этапе build используется образ \texttt{dotnet/sdk:9.0} для компиляции приложения в режиме Release, на этапе runtime используется \texttt{dotnet/aspnet:9.0} для минимизации размера. Приложение включает 9 основных контроллеров (Users, Modules, Progress, Gamification, AiMentor, Notifications, RIMS Sync, Analytics, Reports), запускается на порте 8080. Dockerfile автоматически применяет миграции БД при старте через код в Program.cs (retry logic: 12 попыток с 5-секундными интервалами).

\textbf{Frontend контейнер (Vue.js 3 + Nginx):}

Dockerfile для Onboarding System реализует двухэтапную сборку: build-этап на \texttt{node:20-alpine} выполняет \texttt{npm run build} (Vite собирает статику в dist/), production-этап на \texttt{nginx:stable-alpine} разворачивает статические файлы с поддержкой SPA-маршрутизации (nginx.conf переводит 404 на index.html). Frontend слушает на порте 80 и обслуживает все страницы: Dashboard, Modules, AI Chat, Leaderboard, Analytics.

\textbf{docker-compose для локального тестирования:}

Подготовлен docker-compose.yml для быстрого локального запуска компонентов системы:

\begin{itemize}
	\item \textbf{Frontend-сервис:} собирается из \texttt{onboarding-frontend/Dockerfile}, доступен на http://localhost:80
	\item \textbf{MailHog-сервис:} для перехвата и просмотра email-уведомлений при тестировании (веб-интерфейс на http://localhost:8025)
	\item \textbf{Bridge-сеть:} оба сервиса подключены к сети onboarding\_net для внутренней коммуникации
\end{itemize}

\textbf{Примечание о backend:} В текущей версии compose backend запускается отдельно через \texttt{dotnet run} (Development) или в отдельном контейнере (Production). Это позволяет разработчикам быстро менять код без пересборки образа. Для полнофункционального Production-развёртывания требуется добавить backend-сервис, SQL Server и Redis в compose.

\subsubsection*{Конфигурация для Development и Production}

\textbf{Development окружение (локальная разработка):}

\begin{itemize}
	\item Backend запускается через \texttt{dotnet run} с конфигурацией из \texttt{appsettings.Development.json}
	\item Frontend доступен на http://localhost:5173 через \texttt{npm run dev} (Vite dev server с hot reload)
	\item SQL Server LocalDB или Express используется для локальной БД (создание таблиц через миграции)
	\item MailHog перехватывает все email (вместо реальной SMTP) на http://localhost:8025
	\item CORS настроен для localhost:3000 и localhost:5173 в Program.cs
	\item Логирование на Debug уровне выводит подробную информацию о синхронизации RIMS и ошибках
\end{itemize}

\textbf{Production окружение (развёртывание в организацию):}

\begin{itemize}
	\item Backend собирается в Release-режиме, запускается в Docker контейнере или на выделённом сервере
	\item Frontend собирается в статические файлы (\texttt{npm run build}) и разворачивается на Nginx
	\item Используется корпоративный SQL Server для хранения данных сотрудников, модулей обучения, прогресса, геймификации
	\item Redis (опционально) используется для кеширования Leaderboard, ответов AI Chat, сессий пользователей
	\item Конфигурация задаётся через переменные окружения Docker или appsettings.json в контейнере
	\item RIMS API подключается ежедневно в 02:00 для синхронизации сотрудников и организационной структуры
	\item Email-уведомления отправляются через реальный SMTP (настройки в appsettings.json)
	\item SignalR Hub на /hubs/notifications обеспечивает real-time уведомления в браузере
\end{itemize}

\subsubsection*{Миграции БД и инициализация данных}

\textbf{Entity Framework Core миграции для Onboarding System:}

\begin{itemize}
	\item \textbf{Основные таблицы:} User, Module, TestAttempt, ActionLog (созданы в начальных миграциях)
	\item \textbf{AddGametables (20260517093755):} создаёт таблицы геймификации:
	\begin{itemize}
		\item GamificationProfile --- профиль пользователя с XP, текущий уровень, дата-время последнего достижения
		\item Achievement --- справочник достижений (типы: CompletedModule, FirstTest, LeaderboardTop, ConsecutiveDays)
		\item UserAchievement --- история полученных достижений с датами
	\end{itemize}
	\item \textbf{AddNotifications (20260519125528):} создаёт таблицы для уведомлений:
	\begin{itemize}
		\item Notification --- шаблоны уведомлений (на email, SMS, in-app)
		\item NotificationHistory --- отправленные уведомления с временем доставки
		\item UserNotificationPreference --- настройки пользователя (отписан/подписан от каких типов)
	\end{itemize}
\end{itemize}

\textbf{Автоматическое применение миграций:}

При запуске приложения Program.cs вызывает код, который:

\begin{enumerate}
	\item Проверяет наличие БД (если нет, создаёт её)
	\item Применяет все ожидающие миграции (ApplyPendingMigrations)
	\item Имеет retry-логику: 12 попыток с 5-секундными интервалами на случай временной недоступности SQL Server
	\item Логирует успешное применение миграций для аудита
\end{enumerate}

\textbf{DbInitializer: начальное заполнение данных:}

При первом запуске в Development окружении DbInitializer загружает:

\begin{itemize}
	\item \textbf{Роли:} Admin, Manager, Employee (для контроля прав доступа)
	\item \textbf{Тестовые пользователи:} несколько юзеров с разными должностями и отделами для тестирования
	\item \textbf{Модули обучения:} базовые курсы (Корпоративная культура, IT Безопасность, Инструменты компании)
	\item \textbf{Вопросы и ответы:} тестовые вопросы для модулей (для проверки знаний)
	\item \textbf{Параметры геймификации:} количество XP за разные действия (CompletedModule: 100 XP, FirstTest: 50 XP)
	\item \textbf{Типы уведомлений:} шаблоны email-писем для разных событий (ModuleCompleted, AchievementUnlocked, LeaderboardUpdate)
\end{itemize}

В Production DbInitializer работает аналогично, но данные загружаются из RIMS при первой синхронизации.

\subsubsection*{Синхронизация с RIMS: ключевой компонент внедрения}

Для успешного внедрения Onboarding System в организацию критическое значение имеет интеграция с корпоративной системой управления кадрами RIMS. Все новые сотрудники и их данные поступают из RIMS, обеспечивая единственный источник истины (Single Source of Truth).

\textbf{RimsIntegrationService: реализация синхронизации:}

\begin{itemize}
	\item \textbf{Расписание:} Автоматическая синхронизация запускается ежедневно в 02:00 (не-рабочее время) через Quartz Scheduler в Program.cs
	\item \textbf{Endpoint:} Доступен также ручной запуск через GET \texttt{/api/rims/sync} для администраторов (например, при добавлении срочного новичка)
	\item \textbf{Импортируемые данные:} 
	\begin{itemize}
		\item ФИО сотрудника, email, phone (для уведомлений)
		\item Должность и отдел (для распределения правильных модулей обучения)
		\item Дата приёма на работу (для определения статуса: новичок, активный)
		\item Статус в организации: новичок (требует полного онбординга), активный, уволен
	\end{itemize}
	\item \textbf{Логирование:} Все операции логируются в ActionLog для аудита:
	\begin{itemize}
		\item 23 новых сотрудника успешно импортированы
		\item 5 сотрудников обновлены (изменена должность)
		\item 2 сотрудника помечены как уволены (скрыты из системы)
	\end{itemize}
	\item \textbf{Обработка ошибок:} Если подключение к RIMS не удалось, scheduler повторяет попытку в следующий час (retry logic)
\end{itemize}

\textbf{Конфигурация подключения к RIMS:}

Все параметры указываются в \texttt{appsettings.json}:

\begin{lstlisting}[language=json, caption=Конфигурация RIMS в appsettings.json]
{
  "RimsApi": {
    "BaseUrl": "https://rims.company.local/api",
    "ApiKey": "your-rims-api-key-here",
    "TimeoutSeconds": 30,
    "RetryAttempts": 3,
    "ScheduleTime": "02:00"
  }
}
\end{lstlisting}

\textbf{Рекомендации для Production:}

\begin{itemize}
	\item \textbf{API Key хранение:} Не хранить в исходном коде. Использовать Azure Key Vault, AWS Secrets Manager или переменные окружения
	\item \textbf{TLS сертификат:} RIMS API должен быть доступен по HTTPS с действительным сертификатом
	\item \textbf{Сетевая доступность:} Убедиться, что сервер с Onboarding System может достичь RIMS (может потребоваться выбелить IP в firewall)
	\item \textbf{Тестовое подключение:} Перед production-развёртыванием проверить синхронизацию на тестовых данных RIMS
	\item \textbf{Резервный план:} На случай недоступности RIMS система продолжает работать с существующими пользователями (новички не добавляются, но обучение старых сотрудников не нарушается)
\end{itemize}

\subsubsection*{Конфигурация внешних сервисов для Production}

\textbf{Groq API для AI Chat Assistant:}

Компонент AiChatAssistant использует Groq API для генерации ответов на вопросы сотрудников по модулям обучения:

\begin{itemize}
	\item \textbf{Параметр конфигурации:} \texttt{"GroqApi:ApiKey"} в appsettings.json
	\item \textbf{Использование:} AiMentorService отправляет вопрос сотрудника к Groq и возвращает ответ в браузер
	\item \textbf{Рекомендация:} Протестировать с реальными вопросами по модулям обучения перед внедрением
\end{itemize}

\textbf{Email уведомления (SMTP):}

NotificationService отправляет email в следующих случаях:
\begin{itemize}
	\item Сотрудник назначен на новый модуль обучения (приглашение)
	\item Дедлайн модуля истекает (напоминание)
	\item Достигнуто новое достижение (похвала + результаты на лидерборде)
	\item Manager получает отчёт по прогрессу команды (еженедельно)
\end{itemize}

Параметры SMTP в appsettings.json:

\begin{lstlisting}[language=json, caption=SMTP конфигурация]
{
  "EmailService": {
    "SmtpServer": "mail.company.local",
    "SmtpPort": 587,
    "SenderEmail": "onboarding@company.local",
    "SenderPassword": "****",
    "EnableSSL": true
  }
}
\end{lstlisting}

\textbf{Redis для кеширования (опционально):}

\begin{itemize}
	\item \textbf{Leaderboard:} кеширует топ-100 пользователей на 1 час (иначе каждый запрос считает из БД)
	\item \textbf{AI ответы:} кеширует часто задаваемые вопросы (TTL 24 часа)
	\item \textbf{Сессии:} может использоваться для распределённых сессий (если несколько инстансов backend)
	\item \textbf{Когда требуется:} При более чем 50 одновременных пользователей рекомендуется включить
\end{itemize}

\textbf{SQL Server конфигурация:}

\begin{itemize}
	\item \textbf{Connection string:} В appsettings.json: \texttt{"DefaultConnection"} указывает на SQL Server
	\item \textbf{Версия:} Поддерживается SQL Server 2019 и выше (или Azure SQL Database)
	\item \textbf{Размер БД:} На 1000 сотрудников: примерно 500 MB (таблицы User, Module, TestAttempt, GamificationProfile, Notifications)
	\item \textbf{Резервные копии:} Рекомендуется полная резервная копия ежедневно + транзакционные логи каждый час
\end{itemize}

\subsubsection*{Документация и процесс внедрения}

\textbf{Материалы для администраторов, доступные в репозитории:}

\begin{itemize}
	\item \textbf{OnboardingSystemBackend/Dockerfile:} готов к production build и push в Docker Registry
	\item \textbf{onboarding-frontend/Dockerfile:} готов для статики на Nginx в Production
	\item \textbf{docker-compose.yml:} для локального тестирования (frontend + MailHog). Требует расширения: добавить backend-сервис, SQL Server, Redis
	\item \textbf{OnboardingSystemBackend/Program.cs:} содержит всю DI-конфигурацию: регистрацию контроллеров, БД, CORS, JWT, SignalR, внешних API
	\item \textbf{appsettings.json и appsettings.Development.json:} примеры конфигураций (требуют заполнения реальными значениями для Production)
	\item \textbf{OnboardingSystemBackend/Migrations/:} все миграции БД готовы к применению (ApplyPendingMigrations при старте)
	\item \textbf{OnboardingSystemBackend/Data/DbInitializer.cs:} загрузка тестовых данных и инициализация (в Production данные из RIMS)
\end{itemize}

\textbf{Пошаговый процесс внедрения в организацию:}

\begin{enumerate}
	\item \textbf{Подготовка Infrastructure:}
	\begin{itemize}
		\item Выделить сервер (Windows Server или Linux) для backend
		\item Подготовить SQL Server (версия 2019+ или Azure SQL Database)
		\item Опционально: развернуть Redis для кеширования
		\item Выделить сервер для frontend (или использовать CDN)
	\end{itemize}
	
	\item \textbf{Конфигурация:}
	\begin{itemize}
		\item Заполнить appsettings.json: connection string SQL Server, RIMS API credentials, SMTP, Groq API Key
		\item Настроить DNS записи: onboarding.company.local, rims-api.company.local
		\item Выбелить IP серверов в firewall для RIMS и SMTP
		\item Если используется TLS: установить сертификаты
	\end{itemize}
	
	\item \textbf{Тестирование:}
	\begin{itemize}
		\item Запустить backend: \texttt{dotnet run} --- проверит подключение к SQL Server и применит миграции
		\item Запустить frontend: \texttt{npm run build} + static serve или через Docker
		\item Проверить RIMS синхронизацию: вызвать GET \texttt{/api/rims/sync}, проверить ActionLog
		\item Проверить email: отправить тестовое уведомление на MailHog (http://localhost:8025)
		\item Проверить AI Chat: задать вопрос AiMentorService, получить ответ от Groq
	\end{itemize}
	
	\item \textbf{Миграция данных:}
	\begin{itemize}
		\item DbInitializer загружает тестовые данные при первом запуске (роли, тестовых пользователей, модули)
		\item Первая синхронизация RIMS добавляет реальных сотрудников (статус: newbie)
		\item Manager для каждого нового сотрудника назначает модули обучения
	\end{itemize}
	
	\item \textbf{Запуск в Production:}
	\begin{itemize}
		\item Backend: Docker контейнер с переменными окружения (connection strings, API keys)
		\item Frontend: Docker контейнер Nginx или статика на CDN
		\item Scheduler синхронизации RIMS: запустится автоматически в 02:00
		\item Monitoring: настроить логирование, alerts на ошибки RIMS, SMTP, Groq API
	\end{itemize}
	
	\item \textbf{Post-Launch:}
	\begin{itemize}
		\item Тренинг для Manager: как назначить модули, просмотреть прогресс, создать достижения
		\item Тренинг для HR: как работают модули, как AI Chat помогает новичкам
		\item Тренинг для новичков: как использовать dashboard, модули, AI Chat, leaderboard
	\end{itemize}
\end{enumerate}

\subsubsection*{Текущее состояние готовности к внедрению}

\textbf{Что готово к Production (100\%):}

\begin{enumerate}
	\item \textbf{Backend контейнер:} Dockerfile multi-stage с автоматическими миграциями, готов к Docker push
	\item \textbf{Frontend контейнер:} Nginx multi-stage с SPA-поддержкой, готов к развёртыванию
	\item \textbf{Core функциональность:} 9 контроллеров (Users, Modules, Progress, Gamification, AiMentor, Notifications, RIMS Sync, Analytics, Reports)
	\item \textbf{Геймификация:} Полностью реализована (Achievement, Leaderboard с XP системой)
	\item \textbf{AI Chat Assistant:} Работает с Groq API
	\item \textbf{RIMS интеграция:} RimsIntegrationService с ежедневным scheduler в 02:00
	\item \textbf{Уведомления:} Email-уведомления + SignalR real-time notifications
	\item \textbf{Миграции БД:} Все миграции готовы и автоматически применяются
	\item \textbf{DbInitializer:} Загружает начальные данные и роли
\end{enumerate}

\textbf{Что требует расширения перед Production (но не критично для запуска):}

\begin{enumerate}
	\item \textbf{docker-compose:} Текущая версия неполная (добавить backend, SQL Server, Redis)
	\item \textbf{Документация:} Нет подробного README (требуется написание инструкций для администраторов)
	\item \textbf{E2E тесты:} Структура подготовлена, но сценарии не реализованы (не блокирует Production, но важно для качества)
	\item \textbf{Monitoring:} Нет встроенного мониторинга ошибок (требуется integrация с Prometheus/ELK/Application Insights)
	\item \textbf{Резервные копии:} Нет встроенных процедур backup/restore SQL Server
\end{enumerate}

\textbf{Результат подготовки к внедрению:}

Система \emph{готова к пилотному запуску} в корпоративную сеть организации:

\begin{itemize}
	\item Ежедневная синхронизация с RIMS обеспечит автоматическое добавление новых сотрудников
	\item Модули обучения могут быть назначены Manager сразу после добавления новичка из RIMS
	\item Геймификация и AI Chat Assistant обеспечат engagement новых сотрудников
	\item Email-уведомления и real-time notifications будут информировать о прогрессе обучения
	\item Аналитика (Reports контроллер) предоставит данные для оценки эффективности онбординга
	\item Система может масштабироваться горизонтально (несколько инстансов backend + load balancer)
\end{itemize}

\textbf{Рекомендуемый план запуска:}

\begin{enumerate}
	\item \textbf{Неделя 1-2:} Пилот на 1 отделе (10-20 новичков) с полной проверкой RIMS sync, email, UI
	\item \textbf{Неделя 3-4:} Расширение на 3-4 отдела, сбор feedback от Manager и HR
	\item \textbf{Месяц 2-3:} Полный запуск: все новые сотрудники организации, постоянное улучшение модулей
	\item \textbf{Месяц 4+:} Стабилизация, аналитика эффективности, доработка по feedback
\end{enumerate}

Система готова к внедрению в организации и может обеспечить централизованное управление адаптацией и обучением новых сотрудников с использованием современных подходов (AI, геймификация, аналитика).
